with affine transition morphisms, the canonical map}
\sgaStableTarget{sga6.e13.eq.3-1-1}@@M0@@
\emph{is bijective.}

\emph{Consequently} (EGA IV 8.14.2), \emph{if the Picard functor is represented by an @@M1@@-scheme @@M2@@, then @@M3@@ is locally of finite presentation.}

Indeed, an element @@M4@@ of @@M5@@ is represented by an invertible sheaf @@M6@@ on an @@M7@@ where @@M8@@ is an f.p.p.f. morphism. Put @@M9@@. The image @@M10@@ of @@M11@@ in @@M12@@ is represented by the invertible sheaf @@M13@@ on @@M14@@, the pullback of @@M15@@, where @@M16@@.

If @@M17@@ is the identity element of @@M18@@, there exists an f.p.p.f. morphism @@M19@@ such that the pullback of @@M20@@ to @@M21@@ is @@M22@@. By (EGA IV 8.8.2, 8.10.5 (vi), 11.2.6 (ii), 8.5.2 (i), and 8.5.2.4), there exist a @@M23@@ and an f.p.p.f. morphism @@M24@@, where @@M25@@, such that the pullback of @@M26@@ to @@M27@@ is @@M28@@. The image of @@M29@@ is therefore already the identity element in @@M30@@. Consequently, 3.1.1 is injective.

The surjectivity of 3.1.1 is proved in an analogous way.

\sgaStableTarget{sga6.e13.prop.3-2}\textbf{Proposition 3.2.} \emph{Let @@M31@@ be a noetherian scheme and @@M32@@ a proper @@M33@@-scheme. We suppose that the Picard functor is represented by an @@M34@@-scheme @@M35@@. Then}~:
\begin{enumerate}[label=(\roman*)]
\item \emph{@@M36@@ is locally of finite type over @@M37@@.}
